% 第三部分 B：不依赖结果的章节（由 main_zh.tex 引入）。

\section{设定：只给部分解答打标签的验证器}\label{sec:setting}

第二部分中每个解答都得到验证器的判定。这里每组解答中只有比例 $f=0.25$ 被打标签，即 8 个中随机选 2 个。打标签的比例模拟运行代价高昂的验证器，例如人工评分或缓慢的测试。问题是如何利用另外 75\% 未打标签的解答：
\begin{itemize}
\item \texttt{rlvr\_sub}：把它们排除在外。优势只来自每组中两个打了标签的解答，其余解答没有梯度。
\item \texttt{sub\_kev}：用发布版、冻结的 \kev{} 的 $\pc$（“所给解答的最终答案是否正确？”）作为它们的奖励。
\item \texttt{sub\_kev\_recal}：用 \kev{} 指针头的一个副本给出的 $\pc$ 作为奖励，该副本在迄今见过的所有已标注解答上在线重新拟合。
\item \texttt{sub\_probe}：用一个形状相同、基于相同特征、在相同标签上从头训练的输出头作为奖励。它区分了“\kev{} 发布版输出头的贡献”与“任何拟合这些标签的小输出头都能做到的”。
\end{itemize}
已标注的解答总是得到验证器的奖励。上限参照是第二部分的 \texttt{rlvr}，其中每个解答都有标签。其余一切都沿用第二部分的设置：相同的模型、数据、学习率、批大小、步数和评估（第二部分第 2 节）。仅为记录之用，验证器也检查每个未标注的解答，从而测量每种补全奖励在它所用于的解答上有多准确。

\section{在线重新校准}\label{sec:recal}

\kev{} 的指针头把骨干的隐藏状态变成各选项的 logit $z$，其 $p(\text{yes})$ 为 $\mathrm{softmax}(z/T)_{\text{yes}}$，$T$ 为发布版的温度。\texttt{sub\_kev\_recal} 保持骨干和 \kev{} 的适配器冻结，每一步之后重新拟合输出头的一个 float32 副本 $\theta$。损失为 $p(\text{yes})$ 对验证器标签的二元交叉熵，加上把 $\theta$ 锚定到发布版输出头 $\theta_0$ 的项 $\alpha\lVert\theta-\theta_0\rVert^2$（$\alpha=10^{-3}$）。拟合在最近 8192 个已标注解答的缓冲区上做 50 步 Adam（学习率 $10^{-4}$，批大小 256）。特征（输出头的输入）在 \kev{} 给解答打分时只计算一次，所以重新拟合不需要再跑骨干。骨干上的发布版输出头始终不变。\texttt{sub\_probe} 采用同样的流程，但从输出头投影层的全新初始化开始，$T=1$，不加锚定项。

之所以预期重新校准会起作用，是因为策略的解答在训练中会漂移：变长或变短，风格改变，于是一个在自身训练分布上校准过的评判模型，在这些解答上可能失去校准。已标注的解答直接测量了这种漂移，而在线重新拟合是跟上它的最便宜的办法。
